\vfill\pagebreak
\section{Values known only at run time}
\label{sec:control:runtime_values}

A program decides by asking a question about a value, such as ``is the age at
least 18?'', and the answer chooses what it does next. In the programs so far,
though, every value was written in the source code. If a program contains
\token{let age = 41;}, the question ``is the age at least 18?'' gets the same
answer every time the program runs: yes. Asking it is pointless, since the
program could print its answer directly. Ymir goes further and refuses such a
question outright, as Section~\ref{sec:control:if} shows.

A decision is only useful when its answer can change from one run to the next.
For that, the program needs values that are not written in its source code,
values that it learns only while it runs.

The simplest source of such values is the person running the program. The
function \token{read}, from the module \token{std::io} that already provides
\token{println}, waits for the user to type a line and returns what was typed.

\begin{lstlisting}[style=coloredverbatim, caption=Reading a number typed by the user, label=lst:control:read_age]
use std::io;

fn main() {
  let age = read!i32(ask-> "How old are you? ");
  println("Next year you will be ", age + 1, ".");
}
\end{lstlisting}
\vspace{-10pt}%
\begin{lstlisting}[style=bashVerb, escapechar=@+]
@+\textcolor{teal!80}{alice@dev:\mtilde}@+$ gyc main.yr -o main
@+\textcolor{teal!80}{alice@dev:\mtilde}@+$ ./main
How old are you? 41
Next year you will be 42.
\end{lstlisting}

The call on line~4 has two parts that the book has not presented yet. Until
Section~\ref{sec:functions:read} explains them, treat the whole expression as a
fixed formula:

\begin{itemize}
\item \token{!i32} says which type of value to read. Any other integer type
  works in its place; this chapter reads only \token{i32}.
\item \token{ask-> "..."} gives the text to print before waiting, so the user
  knows what is expected. It can be left out, as in \token{read!i32()}.
\end{itemize}

\token{read} does not check what the user types. If the line does not start
with a number, \token{read!i32} returns \token{0}, and the program carries on
with it. The
programs of this chapter are small enough to live with that.

The second source of values is chance. The function \token{uniform}, from the
module \token{std::rand}, returns a number picked at random between two bounds,
both included. The program below prints a different number, from 1 to 6, each
time it runs.

\begin{lstlisting}[style=coloredverbatim, caption=Rolling a die, label=lst:control:die]
use std::io;
use std::rand;

fn main() {
  let die = uniform(1, 6);
  println("The die shows ", die, ".");
}
\end{lstlisting}

With \token{read} and \token{uniform}, a program can hold values that the
compiler cannot know in advance. The rest of the chapter uses them to make
decisions.
